\section{[\ 附录二\ ]古代汉语常用副词简表}

\noindent
\begin{tabular}{|lp{2em}|p{10em}|p{4em}|p{4em}|}
    \hline
    \multicolumn{2}{|l|}{类别}  & 古代汉语 & 现代汉语 & 备注\\ \hline
    \multicolumn{1}{|l|}{\multirow{15}{*}{\begin{tabular}[c]{@{}l@{}}人\\ 称\\ 代\\ 词\end{tabular}}} & \multirow{5}{*}{\begin{tabular}[c]{@{}l@{}}第\\ 一\\ 人\\ 称\end{tabular}} & 吾、我、余（予） & 我、我们 & 一般自称\\ \cline{3-5} 
    \multicolumn{1}{|l|}{} &  & 某、愚、仆、臣、鄙人、小人、不佞、自称名字 & 我 & 男子谦称\\ \cline{3-5} 
    \multicolumn{1}{|l|}{} &  & 妾、奴、婢子 & 我 & 女子谦称\\ \cline{3-5} 
    \multicolumn{1}{|l|}{} &  & 孤、寡人、孤家、不谷、朕（用于秦始皇之后）& 我 & 帝王谦称\\ \cline{3-5} 
    \multicolumn{1}{|l|}{} &  & “吾”后面加“属、等、侪、辈、曹”等 & 我们这班人 & 复数谦称\\ \cline{2-5} 
    \multicolumn{1}{|l|}{} & \multirow{7}{*}{\begin{tabular}[c]{@{}l@{}}第\\ 二\\ 人\\ 称\end{tabular}} & 汝（女）、尔（而）、若、乃  & 你、你们；你（们）的   & 一般称呼\\ \cline{3-5} 
    \multicolumn{1}{|l|}{} &  & 大人 & 您  & 敬称父母\\ \cline{3-5} 
    \multicolumn{1}{|l|}{} &  & 公、卿、君、子、足下、阁下、夫子、先生、丈人、吾子 & 您  & 敬称一般人\\ \cline{3-5} 
    \multicolumn{1}{|l|}{} &  & 称呼对方的字或官名  & 您 & 敬称一般人\\ \cline{3-5} 
    \multicolumn{1}{|l|}{} &  & 将军、麾下   & 您 & 敬称将军\\ \cline{3-5} 
    \multicolumn{1}{|l|}{} &  & 王、大王、君、上、陛下、殿下   & 您 & 敬称帝王\\ \cline{3-5} 
    \multicolumn{1}{|l|}{} &  & “若、汝、尔”后面加“属、等、辈、曹”等  & 你们这班人   & 复数卑称\\
    \hline
\end{tabular}
\newpage

\begin{flushright}
    （续表）
\end{flushright}
\noindent
\begin{tabular}{|lp{2em}|p{8em}|p{6em}|p{4em}|}
    \hline
    \multicolumn{2}{|l|}{类别}  & 古代汉语 & 现代汉语 & 备注\\ \hline
    \multicolumn{1}{|l|}{\multirow{4}{*}{\begin{tabular}[c]{@{}l@{}}人\\ 称\\ 代\\ 词\end{tabular}}} & \multirow{3}{*}{\begin{tabular}[c]{@{}l@{}}第\\ 三\\ 人\\ 称\end{tabular}} & 其、厥   & 他（它）的，他（它）们的 & 用作定语\\ \cline{3-5} 
    \multicolumn{1}{|l|}{} &  & 之  & 他、他们 & 用作宾语\\ \cline{3-5} 
    \multicolumn{1}{|l|}{} &  & 彼  & 他、他们 & 用作主语\\ \hline

    \multicolumn{1}{|l|}{\multirow{8}{*}{\begin{tabular}[c]{@{}l@{}}指\\ 示\\ 代\\ 词\end{tabular}}} & 远指 & 彼、其、夫（用在句子中） & 那、那个、那些、那样    & “彼”可指人\\ \cline{2-5} 
    \multicolumn{1}{|l|}{}  & \multirow{3}{*}{近指} & 是、此、兹、斯 & 这、这个、这些、这样、这儿 & “是”\hspace{1em}“此”可指人\\ \cline{3-5} 
    \multicolumn{1}{|l|}{}  &   & 然、尔 & 这样、如此 & 一般用作谓语\\ \cline{3-5} 
    \multicolumn{1}{|l|}{}  &   & 之  & 这、这样、这些  & 一般用作定语\\ \cline{2-5} 
    \multicolumn{1}{|l|}{}  & 旁指  & 他   & 别的、其他的  & 与现代用法不同\\ \cline{2-5} 
    \multicolumn{1}{|l|}{}  & \multirow{3}{*}{无定} & 或   & 有的、有人、有的人     & 不同于“或者”\\ \cline{3-5} 
    \multicolumn{1}{|l|}{}  &   & 某  & 某某人  & 不限于代“我”\\ \cline{3-5} 
    \multicolumn{1}{|l|}{}  &   & 莫（作否定副词用时，相当于“勿”，“不”） & 没有谁、没有什么东西    &                   \\ \hline
\end{tabular}
\newpage

\begin{flushright}
	（续表）
\end{flushright}
\noindent
\begin{tabular}{|lp{2em}|p{8em}|p{6em}|p{4em}|}
    \hline
    \multicolumn{2}{|l|}{类别} & 古代汉语 & 现代汉语 & 备注\\ \hline
    \multicolumn{1}{|l|}{\multirow{3}{*}{\begin{tabular}[c]{@{}l@{}}疑\\ 问\\ 代\\ 词\end{tabular}}} & 问人  & 谁、孰   & 谁、什么  & “孰”兼问事物  \\ \cline{2-5} 
    \multicolumn{1}{|l|}{}  & 问事物   & 何、奚、胡、曷    & 怎么、为什么  & “何”兼有“什么、哪、哪里”的意思\\ \cline{2-5} 
    \multicolumn{1}{|l|}{}  & 问处所或反问  & 安、焉、恶（wū）、乌  & 怎么、哪里  & 经常作状语，有时作宾语\\ \hline
\end{tabular}